\newcommand{\bMedio}{1.253314}
\def\bInversa{0.3989422804014327}
\def\aIntercept{1.177410} % $\ln(4)$
\def\sAB{0.8857913124365199} % $\frac{a}{b-1}$
\def\b{2.50662}
% \section{Gaussiana}

\begin{frame}{Monty Python}{\cite{Marsaglia-MontyPython}}
% 
  \begin{tikzpicture}
    % b -----> sqrt(2*pi)=2.50662
    % b/2 ---> 1.253314
    % 1/b ---> 0.39894
    \begin{axis}[
        domain=0:4, % Rango de los valores en el eje x
        samples=100, % Cantidad de puntos a graficar
        axis lines=middle, % Coloca los ejes en el medio
        xlabel=$x$, % Etiqueta del eje x
        ylabel={$f(x)$}, % Etiqueta del eje y
        xmin=0, xmax=3.5, % Limites del eje x
        ymin=-0.2, ymax=0.9, % Limites del eje y
        xtick={0,1,2,3}, % Puntos de referencia en el eje x
        ytick={0,0.1,0.2,...,0.8}, % Puntos de referencia en el eje y
        domain=0:4, % Dominio de la función
        samples=200, % Resolución de la curva
        smooth % Curva suave
      ]
      \addplot[blue, thick] {2/(sqrt(2*pi)) * exp(-0.5*x^2)};
% 
% 
      % Área bajo la curva (relleno)
      \addplot [
        fill=blue!30, % Color de relleno (30% de azul)
        domain=0:4,
      ]
      {2*(1/sqrt(2*pi)) * exp(-0.5*x^2)} \closedcycle; % Relleno bajo la curva  
% 
      % Segunda gaussiana rotada 45 grados
      \begin{scope}[rotate around={90:(\aIntercept,0.3989422)}] % Rotar alrededor del origen 45 grados
        % Curva Gaussiana rotada
        \addplot [
          thick, % Grosor de la línea
          green!30, % Color de la línea
          domain=0:\aIntercept, % Dominio en el eje x
        ]
        {2 * (1/sqrt(2*pi)) * exp(-0.5*x^2)}; % Ecuación de la curva Gaussiana multiplicada por 2
% 
        % Área bajo la curva rotada (relleno)
        \addplot [
          fill=green!30, % Color de relleno (30% de verde)
          domain=0:\aIntercept,
        ]
        {2 * (1/sqrt(2*pi)) * exp(-0.5*x^2)} \closedcycle;
      \end{scope}
% 
      % Dibujar el rectángulo de ancho sqrt(2*pi) y altura 1/ancho
      \draw[fill=red!30, draw=red, opacity=0.5]
      (0, 0) rectangle ({sqrt(2*pi)}, \bInversa);
      % (0, 0) rectangle ({sqrt(2*pi)}, 0.363792);
% 
      %       % Derecha del rectángulo
      %       \addplot [
      %         fill=cyan!50, % Color amarillo para las áreas fuera del rectángulo
      %         domain={sqrt(2*pi)/2}:4,
      %       ]
      %       {2 * (1/sqrt(2*pi)) * exp(-0.5*x^2)} \closedcycle;
      %       % Cola
      %       \addplot [
      %         fill=green!50, % Color amarillo para las áreas fuera del rectángulo
      %         domain={sqrt(2*pi)}:4,
      %       ]
      %       {2 * (1/sqrt(2*pi)) * exp(-0.5*x^2)} \closedcycle;
      % %   Rectángulo corrector (blanco)
      %       \draw[fill=white, draw=white] 
      %       (\bMedio, 0.405) rectangle (3,0.8);
      \draw [-Triangle,red] (axis cs: \b,0) -- (axis cs: \b,-0.05)node[below]{\scriptsize $b$};
      \draw [-Triangle,red] (axis cs: \aIntercept,0) -- (axis cs: \aIntercept,-0.05)node[below]{\scriptsize $a$};
    \end{axis}
  \end{tikzpicture}
% 
\end{frame}
% 
\begin{frame}{Traslados}{Transformaciones}

%   \begin{tikzpicture}
%     % b -----> sqrt(2*pi)=\b
%     % b/2 ---> 1.253314
%     % 1/b ---> 0.39894
%     \def\yMax{1.0}
%     \begin{axis}[
%         domain=0:4, % Rango de los valores en el eje x
%         samples=100, % Cantidad de puntos a graficar
%         axis lines=middle, % Coloca los ejes en el medio
%         xlabel=$x$, % Etiqueta del eje x
%         ylabel={$f(x)$}, % Etiqueta del eje y
%         xmin=0, xmax=3.5, % Limites del eje x
%         ymin=-0.2, ymax=\yMax, % Limites del eje y
%         xtick={0,1,2,3}, % Puntos de referencia en el eje x
%         ytick={0,0.1,0.2,...,\yMax}, % Puntos de referencia en el eje y
%         domain=0:4, % Dominio de la función
%         samples=200, % Resolución de la curva
%         smooth % Curva suave
%       ]
%       \addplot[blue, thick] {2/(sqrt(2*pi)) * exp(-0.5*x^2)};


%       % % ! Área bajo la curva (relleno)
%       \addplot [
%         fill=blue!30, % Color de relleno (30% de azul)
%         domain=0:4,
%       ]
%       {2*(1/sqrt(2*pi)) * exp(-0.5*x^2)} \closedcycle; % Relleno bajo la curva
%       % % ! Dibujar el rectángulo de ancho sqrt(2*pi) y altura 1/ancho
%       \draw[fill=orange!30, draw=orange, opacity=0.5]
%       (0, 0) rectangle ({sqrt(2*pi)}, \bInversa);

%       \begin{scope}[shift={(0.58870,-\bInversa)}] % Trasladar en ejes x,y
%         % % ! Área bajo la curva rotada (relleno)
%         \addplot [
%           fill=blue!70, % Color de relleno (30% de verde)
%           domain=0:\aIntercept
%         ]
%         {(2/sqrt(2*pi)) * exp(-0.5*(x)^2)} \closedcycle;
%       \end{scope}
%       \begin{scope}[shift={(0.58870,-0.20)}] % Trasladar en ejes x,y
%         % % ! Área bajo la curva rotada (relleno)
%         \addplot [
%           fill=green!70, % Color de relleno (30% de verde)
%           domain=0:\aIntercept
%         ]
%         {(2/sqrt(2*pi)) * exp(-0.5*(x)^2)} \closedcycle;
%       \end{scope}

%       \draw[-,red,dashed] (0,0.8)--(2,0.8);
%     \end{axis}
%   \end{tikzpicture}
\end{frame}
% 
% 
\begin{frame}{Traslados}{Transformaciones}
% 
  \begin{tikzpicture}
    % b -----> sqrt(2*pi)=\b
    % b/2 ---> 1.253314
    % 1/b ---> 0.39894
    \def\yMax{1.0}
    \begin{axis}[
        domain=0:4, % Rango de los valores en el eje x
        samples=100, % Cantidad de puntos a graficar
        axis lines=middle, % Coloca los ejes en el medio
        xlabel=$x$, % Etiqueta del eje x
        ylabel={$f(x)$}, % Etiqueta del eje y
        xmin=0, xmax=3.5, % Limites del eje x
        ymin=-0.2, ymax=\yMax, % Limites del eje y
        xtick={0,1,2,3}, % Puntos de referencia en el eje x
        ytick={0,0.1,0.2,...,\yMax}, % Puntos de referencia en el eje y
        domain=0:4, % Dominio de la función
        samples=200, % Resolución de la curva
        smooth % Curva suave
      ]
      % % ! Área bajo la curva (relleno)
      \addplot [
        fill=blue!30, % Color de relleno (30% de azul)
        draw=none,
        domain=0:4,
      ]
      {2*(1/sqrt(2*pi)) * exp(-0.5*x^2)} \closedcycle; % Relleno bajo la curva
      % % ! Dibujar el rectángulo de ancho sqrt(2*pi) y altura 1/ancho
      \draw[fill=orange!30, draw=none, opacity=0.5]
      (0, 0) rectangle ({sqrt(2*pi)}, \bInversa);
% 
      % \begin{scope}[shift={(0.58870,-\bInversa)}] % Trasladar en ejes x,y
      %   % % ! Área bajo la curva rotada (relleno)
      %   \addplot [
      %     fill=blue!70, % Color de relleno (30% de verde)
      %     domain=0:\aIntercept
      %   ]
      %   {(2/sqrt(2*pi)) * exp(-0.5*(x)^2)} \closedcycle;
      % \end{scope}
      % % ! Rotada respecto al eje x
      \begin{scope}[shift={(0.58870,-0.20)}, rotate around={90:(0.958,0.3)}] % Trasladar en ejes x,y
        \addplot [
          fill=green!40, % Color de relleno (30% de verde)
          draw=none,
          domain=0:\aIntercept
        ]
        % {(2/sqrt(2*pi)) * exp(-0.5*(x)^2)} \closedcycle;
        {(2/sqrt(2*pi)) * exp(-0.4428956*(x)^2)} \closedcycle; % ! Estiramiento horizontal
      \end{scope}
      \filldraw (axis cs: \aIntercept,\bInversa) circle (2pt);
% 
      \draw[-,red,dashed] (0,0.8)--(2,0.8);
      \draw [-Triangle,red] (axis cs: \b,0) -- (axis cs: \b,-0.05)node[below]{\scriptsize $b$};
      \draw [-Triangle,red] (axis cs: \aIntercept,0) -- (axis cs: \aIntercept,-0.05)node[below]{\scriptsize $a$};
    \end{axis}
  \end{tikzpicture}
\end{frame}
% \def\b{2.50662} % \def\bMedio{1.253314}
% % \def\bInversa{0.3989422804014327}
% % \def\aIntercept{1.177410} % $\ln(4)$
% % \def\sAB{0.8857913124365199} % $\frac{a}{b-1}$


% \section{von Mises}
% \begin{frame}{Traslados}{Transformaciones}
%   \def\BesselK{1.5}
%   \def\BesselI{0.5118}
%   \begin{tikzpicture}
%     % b -----> sqrt(2*pi)=\b
%     % b/2 ---> 1.253314
%     % 1/b ---> 0.39894 
%     \def\yMax{1.0}
%     \begin{axis}[
%         domain=0:4, % Rango de los valores en el eje x
%         samples=100, % Cantidad de puntos a graficar
%         axis lines=middle, % Coloca los ejes en el medio
%         xlabel=$x$, % Etiqueta del eje x
%         ylabel={$f(x)$}, % Etiqueta del eje y
%         xmin=0, xmax=3.5, % Limites del eje x
%         ymin=-0.2, ymax=\yMax, % Limites del eje y
%         xtick={0,1,2,3}, % Puntos de referencia en el eje x
%         ytick={0,0.1,0.2,...,\yMax}, % Puntos de referencia en el eje y
%         domain=0:4, % Dominio de la función
%         samples=200, % Resolución de la curva
%         smooth % Curva suave
%       ]
%       % % ! Área bajo la curva (relleno)
%       \addplot [
%         fill=orange!30, % Color de relleno (30% de azul)
%         draw=none,
%         domain=0:3.1415,
%       ]
%       {(1/\BesselI*pi) * exp(\BesselK*cos(deg(x)))} \closedcycle; % Relleno bajo la curva
%       % % ! Dibujar el rectángulo de ancho sqrt(2*pi) y altura 1/ancho
%       \draw[fill=orange!30, draw=none, opacity=0.5]
%       (0, 0) rectangle ({sqrt(2*pi)}, \bInversa);
%       % % % ! Rotada respecto al eje x
%       % \begin{scope}[shift={(0.58870,-0.20)}, rotate around={90:(0.958,0.3)}] % Trasladar en ejes x,y
%       %   \addplot [
%       %     fill=green!40, % Color de relleno (30% de verde)
%       %     draw=none,
%       %     domain=0:\aIntercept
%       %   ]
%       %   % {(2/sqrt(2*pi)) * exp(-0.5*(x)^2)} \closedcycle;
%       %   {(2/sqrt(2*pi)) * exp(-0.4428956*(x)^2)} \closedcycle; % ! Estiramiento horizontal
%       % \end{scope}
%       \filldraw (axis cs: \aIntercept,\bInversa) circle (2pt);
%       \draw [-Triangle,red] (axis cs: \b,0) -- (axis cs: \b,-0.05)node[below]{\scriptsize $b$};
%       \draw [-Triangle,red] (axis cs: \aIntercept,0) -- (axis cs: \aIntercept,-0.05)node[below]{\scriptsize $a$};
%     \end{axis}
%     \node (A) at (8.2,5){$f(x) = \left(\frac{1}{\pi I_0(z)}\right)e^{z\,\cos(x)}$};
%   \end{tikzpicture}
% \end{frame}

% \section{Monty Python}

% \subsection{Algoritmo Monty Python}
% \begin{frame}{Monty Python}{Algoritmo}
%   \begin{itemize}
%     \item Cita de \parencite{Marsaglia-MontyPython}\footfullcite{Marsaglia-MontyPython}
%   \end{itemize}
% \end{frame}

